\section{Normal maps between ordinals}

\import{cardinal/cofinality.tex}

A normal map between ordinals is strictly increasing and continuous.  A
normal cofinal map additionally has cofinal range, and a normal cofinal
cardinal sequence additionally has cardinal values.  Index domination and
infinite-cardinal-valuedness are separate refinements. The component
increasingness, continuity, value, and domination predicates are deliberately
raw bounded conditions; the full normal-map predicates supply the ordinal and
function typing.

\begin{definition}\label{ordinal_map}
    $\ordinalMap{c}{\lambda}{\kappa}$ iff
        $\lambda$ and $\kappa$ are ordinals and
        $c\in\funs{\lambda}{\kappa}$.
\end{definition}

\begin{definition}\label{normal_map_increasing}
    $\normalMapIncreasing{c}{\lambda}$ iff
        for all $j,k$ such that $j\in k\in\lambda$ we have $c(j)\in c(k)$.
\end{definition}

\begin{definition}\label{normal_map_continuous}
    $\normalMapContinuous{c}{\lambda}$ iff
        for every limit ordinal $l\in\lambda$ we have
            $c(l)=\unions{\ran{\restrl{c}{l}}}$.
\end{definition}

\begin{definition}\label{normal_map}
    $\normalMap{c}{\lambda}{\kappa}$ iff
        $\ordinalMap{c}{\lambda}{\kappa}$ and
        $\normalMapIncreasing{c}{\lambda}$ and
        $\normalMapContinuous{c}{\lambda}$.
\end{definition}

\begin{definition}\label{normal_cofinal_map}
    $\normalCofinalMap{c}{\lambda}{\kappa}$ iff
        $\normalMap{c}{\lambda}{\kappa}$ and
        $\cofinalMap{c}{\lambda}{\kappa}$.
\end{definition}

\begin{definition}\label{cardinal_sequence_values}
    $\cardinalSequenceValues{c}{\lambda}$ iff
        for every $i\in\lambda$ we have $c(i)$ is a cardinal.
\end{definition}

\begin{definition}\label{infinite_cardinal_sequence_values}
    $\infiniteCardinalSequenceValues{c}{\lambda}$ iff
        for every $i\in\lambda$ we have $c(i)$ is an infinite cardinal.
\end{definition}

\begin{definition}\label{normal_cofinal_cardinal_sequence}
    $\normalCofinalCardinalSequence{c}{\lambda}{\kappa}$ iff
        $\normalCofinalMap{c}{\lambda}{\kappa}$ and
        $\cardinalSequenceValues{c}{\lambda}$.
\end{definition}

\begin{definition}\label{index_dominating}
    $c$ is index-dominating on $\lambda$ iff
        for every $i\in\lambda$ we have $\card{i}\in c(i)$.
\end{definition}

\begin{proposition}\label{normal_map_type}
    Suppose $\normalMap{c}{\lambda}{\kappa}$.
    Then $\lambda$ and $\kappa$ are ordinals and
        $c\in\funs{\lambda}{\kappa}$.
\end{proposition}
\begin{proof}
    Follows by \auto{normal_map,ordinal_map}.
\end{proof}

\begin{proposition}\label{normal_map_increasing_elim}
    Suppose $\normalMap{c}{\lambda}{\kappa}$.
    Let $k\in\lambda$.
    Let $j\in k$.
    Then $c(j)\in c(k)$.
\end{proposition}
\begin{proof}
    $\normalMapIncreasing{c}{\lambda}$ by \obv{normal_map}.
    Follows by \obv{normal_map_increasing}.
\end{proof}

\begin{proposition}\label{normal_map_continuous_elim}
    Suppose $\normalMap{c}{\lambda}{\kappa}$.
    Let $l\in\lambda$.
    Suppose $l$ is a limit ordinal.
    Then $c(l)=\unions{\ran{\restrl{c}{l}}}$.
\end{proposition}
\begin{proof}
    $\normalMapContinuous{c}{\lambda}$ by \obv{normal_map}.
    Follows by \obv{normal_map_continuous}.
\end{proof}

\begin{proposition}\label{normal_cofinal_map_is_normal}
    Suppose $\normalCofinalMap{c}{\lambda}{\kappa}$.
    Then $\normalMap{c}{\lambda}{\kappa}$.
\end{proposition}
\begin{proof}
    Follows by \obv{normal_cofinal_map}.
\end{proof}

\begin{proposition}\label{normal_cofinal_map_is_cofinal}
    Suppose $\normalCofinalMap{c}{\lambda}{\kappa}$.
    Then $\cofinalMap{c}{\lambda}{\kappa}$.
\end{proposition}
\begin{proof}
    Follows by \obv{normal_cofinal_map}.
\end{proof}

\begin{proposition}\label{normal_cofinal_map_type}
    Suppose $\normalCofinalMap{c}{\lambda}{\kappa}$.
    Then $\lambda$ and $\kappa$ are ordinals and
        $c\in\funs{\lambda}{\kappa}$.
\end{proposition}
\begin{proof}
    $\normalMap{c}{\lambda}{\kappa}$
        by \obv{normal_cofinal_map_is_normal}.
    Follows by \obv{normal_map_type}.
\end{proof}

\begin{proposition}\label{normal_cofinal_map_increasing}
    Suppose $\normalCofinalMap{c}{\lambda}{\kappa}$.
    Let $k\in\lambda$.
    Let $j\in k$.
    Then $c(j)\in c(k)$.
\end{proposition}
\begin{proof}
    $\normalMap{c}{\lambda}{\kappa}$
        by \obv{normal_cofinal_map_is_normal}.
    Follows by \obv{normal_map_increasing_elim}.
\end{proof}

\begin{proposition}\label{normal_cofinal_map_continuous}
    Suppose $\normalCofinalMap{c}{\lambda}{\kappa}$.
    Let $l\in\lambda$.
    Suppose $l$ is a limit ordinal.
    Then $c(l)=\unions{\ran{\restrl{c}{l}}}$.
\end{proposition}
\begin{proof}
    $\normalMap{c}{\lambda}{\kappa}$
        by \obv{normal_cofinal_map_is_normal}.
    Follows by \obv{normal_map_continuous_elim}.
\end{proof}

\begin{proposition}\label{normal_cofinal_cardinal_sequence_is_normal_cofinal}
    Suppose $\normalCofinalCardinalSequence{c}{\lambda}{\kappa}$.
    Then $\normalCofinalMap{c}{\lambda}{\kappa}$.
\end{proposition}
\begin{proof}
    Follows by \obv{normal_cofinal_cardinal_sequence}.
\end{proof}

\begin{proposition}\label{normal_cofinal_cardinal_sequence_type}
    Suppose $\normalCofinalCardinalSequence{c}{\lambda}{\kappa}$.
    Then $\lambda$ and $\kappa$ are ordinals and
        $c\in\funs{\lambda}{\kappa}$.
\end{proposition}
\begin{proof}
    $\normalCofinalMap{c}{\lambda}{\kappa}$
        by \obv{normal_cofinal_cardinal_sequence}.
    Hence $\normalMap{c}{\lambda}{\kappa}$
        by \obv{normal_cofinal_map_is_normal}.
    Follows by \obv{normal_map_type}.
\end{proof}

\begin{proposition}\label{normal_cofinal_cardinal_sequence_is_cofinal}
    Suppose $\normalCofinalCardinalSequence{c}{\lambda}{\kappa}$.
    Then $\cofinalMap{c}{\lambda}{\kappa}$.
\end{proposition}
\begin{proof}
    $\normalCofinalMap{c}{\lambda}{\kappa}$
        by \obv{normal_cofinal_cardinal_sequence}.
    Follows by \obv{normal_cofinal_map_is_cofinal}.
\end{proof}

\begin{proposition}\label{normal_cofinal_cardinal_sequence_value}
    Suppose $\normalCofinalCardinalSequence{c}{\lambda}{\kappa}$.
    Let $i\in\lambda$.
    Then $c(i)$ is a cardinal.
\end{proposition}
\begin{proof}
    $\cardinalSequenceValues{c}{\lambda}$
        by \obv{normal_cofinal_cardinal_sequence}.
    Follows by \obv{cardinal_sequence_values}.
\end{proof}

\begin{proposition}\label{normal_cofinal_cardinal_sequence_increasing}
    Suppose $\normalCofinalCardinalSequence{c}{\lambda}{\kappa}$.
    Let $k\in\lambda$.
    Let $j\in k$.
    Then $c(j)\in c(k)$.
\end{proposition}
\begin{proof}
    $\normalCofinalMap{c}{\lambda}{\kappa}$
        by \obv{normal_cofinal_cardinal_sequence}.
    Hence $\normalMap{c}{\lambda}{\kappa}$
        by \obv{normal_cofinal_map_is_normal}.
    Follows by \obv{normal_map_increasing_elim}.
\end{proof}

\begin{proposition}\label{normal_cofinal_cardinal_sequence_continuous}
    Suppose $\normalCofinalCardinalSequence{c}{\lambda}{\kappa}$.
    Let $l\in\lambda$.
    Suppose $l$ is a limit ordinal.
    Then $c(l)=\unions{\ran{\restrl{c}{l}}}$.
\end{proposition}
\begin{proof}
    $\normalCofinalMap{c}{\lambda}{\kappa}$
        by \obv{normal_cofinal_cardinal_sequence}.
    Hence $\normalMap{c}{\lambda}{\kappa}$
        by \obv{normal_cofinal_map_is_normal}.
    Follows by \obv{normal_map_continuous_elim}.
\end{proof}

\begin{proposition}\label{index_dominating_value}
    Suppose $c$ is index-dominating on $\lambda$.
    Let $i\in\lambda$.
    Then $\card{i}\in c(i)$.
\end{proposition}
\begin{proof}
    Follows by \obv{index_dominating}.
\end{proof}

\begin{proposition}\label{index_dominating_normal_map_value_contains_zero}
    Suppose $\normalMap{c}{\lambda}{\kappa}$.
    Suppose $c$ is index-dominating on $\lambda$.
    Let $i\in\lambda$.
    Then $\emptyset\in c(i)$.
\end{proposition}
\begin{proof}
    $\card{i}\in c(i)$ by \obv{index_dominating_value}.
    Hence $c(i)$ is inhabited.
    Hence $c(i)\neq\emptyset$ by \auto{nonempty_inhabited}.
    $c\in\funs{\lambda}{\kappa}$ and $\kappa$ is an ordinal
        by \obv{normal_map_type}.
    Hence $c(i)\in\kappa$ by \obv{funs_type_apply}.
    Hence $c(i)$ is an ordinal by \obv{elem_of_ordinal_is_ordinal}.
    Follows by \obv{ordinal_empty_or_emptyset_elem}.
\end{proof}

\begin{proposition}\label{infinite_cardinal_sequence_value}
    Suppose $\infiniteCardinalSequenceValues{c}{\lambda}$.
    Let $i\in\lambda$.
    Then $c(i)$ is an infinite cardinal.
\end{proposition}
\begin{proof}
    Follows by \obv{infinite_cardinal_sequence_values}.
\end{proof}

\begin{proposition}\label{infinite_cardinal_sequence_values_are_cardinal_values}
    Suppose $\infiniteCardinalSequenceValues{c}{\lambda}$.
    Then $\cardinalSequenceValues{c}{\lambda}$.
\end{proposition}
\begin{proof}
    For every $i\in\lambda$ we have $c(i)$ is an infinite cardinal
        by \obv{infinite_cardinal_sequence_values}.
    Hence for every $i\in\lambda$ we have $c(i)$ is a cardinal.
    Follows by \obv{cardinal_sequence_values}.
\end{proof}
